\documentclass[11pt]{article}
\usepackage[utf8]{inputenc}
\usepackage{geometry}
\geometry{a4paper, margin=1in}
\usepackage{enumitem}
\usepackage{hyperref}
\usepackage{fancyhdr}
\usepackage{graphicx}
\usepackage{listings}
\usepackage{xcolor}
\usepackage{titlesec}

% Formatting code blocks
\lstset{
    language=C,
    basicstyle=\footnotesize\ttfamily,
    backgroundcolor=\color{gray!10},
    frame=single,
    rulecolor=\color{gray},
    showstringspaces=false,
    tabsize=4,
    captionpos=b,
    breaklines=true,
    keywordstyle=\color{blue}\bfseries,
    commentstyle=\color{green!60!black}\itshape,
    stringstyle=\color{red},
    numbers=left,
    numberstyle=\tiny\color{gray},
    numbersep=5pt,
}

% Header/Footer
\pagestyle{fancy}
\fancyhf{}
\rhead{CS412 Fuzzing Lab}
\lhead{Fuzzing Lab Report}
\rfoot{Page \thepage}

% Title formatting
\titleformat{\section}{\normalfont\Large\bfseries}{\thesection}{1em}{}
\titleformat{\subsection}{\normalfont\large\bfseries}{\thesubsection}{1em}{}

% Title Information
\title{\textbf{CS412 Fuzzing Lab Report}}
\date{April 2025}

\begin{document}

    \maketitle

    \begin{abstract}
        This is a simplified example report to demonstrate expectations for the fuzzing lab project. The actual report should include more detailed technical discussion and advanced improvements.
    \end{abstract}

    \section{Introduction}
    Students:
    \begin{itemize}
        \item Emile Cornamusaz, 340827
        \item Yann Savioz, 326801
    \end{itemize}
    For the fuzzing lab we chose the \texttt{libpng} project \url{https://github.com/pnggroup/libpng}, a widely used open-source C library for reading and writing PNG image files. \\
    Forked libpng project repository : \url{https://github.com/ecornamu/libpng} \\
    Forked oss-fuzz project repository : \url{https://github.com/ecornamu/oss-fuzz}

    \section{Part 1}
    Libpng has only one harness (libpng\_read\_fuzzer).
    \subsection{Build and run the harness}
    Once the forked repositories are cloned, we proceed as follows to run the drivers for 4 hours with default corpus (commands are ran in the oss-fuzz directory)
    \begin{lstlisting}
python3 infra/helper.py build_image libpng
python3 infra/helper.py build_fuzzers libpng ../libpng --clean
mkdir libpng_with_corpus
timeout 4h python3 infra/helper.py run_fuzzer libpng libpng_read_fuzzer --corpus-dir libpng_with_corpus
    \end{lstlisting}
    To run the harness with an empty corpus, we need to modify the build script located in \texttt{libpng/contrib/oss-fuzz}. The diff is located at \texttt{submission/part1/project.diff} and is as follows :
    \begin{lstlisting}
diff --git a/contrib/oss-fuzz/build.sh b/contrib/oss-fuzz/build.sh
index 7b8f02639..c7d0fce84 100755
--- a/contrib/oss-fuzz/build.sh
+++ b/contrib/oss-fuzz/build.sh
@@ -43,8 +43,8 @@ $CXX $CXXFLAGS -std=c++11 -I. \
      -lFuzzingEngine .libs/libpng16.a -lz
 
 # add seed corpus.
-find $SRC/libpng -name "*.png" | grep -v crashers | \
-     xargs zip $OUT/libpng_read_fuzzer_seed_corpus.zip
+# find $SRC/libpng -name "*.png" | grep -v crashers | \
+#      xargs zip $OUT/libpng_read_fuzzer_seed_corpus.zip
 
 cp $SRC/libpng/contrib/oss-fuzz/*.dict \
      $SRC/libpng/contrib/oss-fuzz/*.options $OUT/

    \end{lstlisting}
    We commented the line that adds the seed corpus to the harness.
    Then we run the following commands to execute the harness without seeds for 4 hours.
    \begin{lstlisting}
cd libpng
git apply ../project.diff
cd ../oss-fuzz
python3 infra/helper.py build_image libpng
python3 infra/helper.py build_fuzzers libpng ../libpng --clean
mkdir libpng_without_corpus
timeout 5m python3 infra/helper.py run_fuzzer libpng libpng_read_fuzzer --corpus-dir libpng_without_corpus
    \end{lstlisting}
    Both scripts are available at \texttt{submission/part1/run.w\_corpus.sh} and \texttt{submission/part1/run.w\_o\_corpus.sh} respectively.
    \subsection{Coverage report}
    The coverage report for the fuzzing with and without corpus are available in \texttt{submission/part1/report/w\_corpus/} and \texttt{submission/part1/report/w\_o\_corpus/} respectively. \\ \\
    Here is the coverage summary for both runs, illustrating the measurable impact of using a seed corpus during fuzzing of the \texttt{libpng} library. All key coverage metrics show improvements when a corpus is included. The most notable gain is in function coverage, which increased by 2\%, indicating that more logical branches and internal operations were successfully exercised. \newline
    \begin{table}[h!]
        \centering
        \begin{tabular}{|l|c|c|c|}
            \hline
            \textbf{Metric} & \textbf{With Corpus} & \textbf{Without Corpus} & \textbf{Difference} \\
            \hline
            \textbf{Functions} & 50.5\% & 48.5\%  & \textbf{+2.0\%} \\
            \textbf{Lines} & 41.6\%  & 40.8\%  & \textbf{+0.8\%} \\
            \textbf{Regions} & 31.4\%  & 30.8\%  & \textbf{+0.6\%} \\
            \hline
        \end{tabular}
        \caption{Overall Coverage Summary: Fuzzing libpng with vs. without seed corpus}
    \end{table}

    First we are going to take a look at key code areas only covered with seed corpus in \texttt{png.c}

    \subsubsection{Zlib Decompression Error Handling}
    \begin{lstlisting}
switch(ret) {                       // line 1004
...
case Z_OK:                        // line 1007
png_ptr->zstream.msg = PNGZ_MSG_CAST("unexpected zlib return code");  // line 1008
break;                           // line 1009
    \end{lstlisting}

    These lines occur during decompression of image data, which is only triggered when valid image data passes multiple layers of chunk and header validation. They are thus hardly achievable without a seed corpus providing valid data.

    \subsubsection{Byte Decoding and Character Mapping}
    \begin{lstlisting}
byte &= 0xff;                    // line 1499
if (byte >= 32 && byte <= 126)  // line 1500
return (char)byte;              // line 1501
    \end{lstlisting}

    This snippet converts raw bytes to readable characters for reporting or debugging. It is invoked during processing of specific chunks that appear in real PNGs, such as metadata tags, and is not reached with random bytes.

    \subsubsection{ICC Profile Handling Logic}
    \begin{lstlisting}
name[1] = png_icc_tag_char(tag >> 24);                 // line 1510
return it == 32 || (it >= 48 && it <= 57) ...          // line 1520
if (is_ICC_signature(value) != 0)                      // line 1543
if (profile_length > png_chunk_max(png_ptr))           // line 1603
    \end{lstlisting}

    ICC (International color consortium) profile validation requires the presence of an iCCP chunk, which contains profile length, tag structure, and character validation. Those requirements are hard to meet with a random input.

    \subsubsection{Structured Error Message Generation}
    \begin{lstlisting}
pos = png_safecat(message, ..., "profile '");         // line 1540
if (is_ICC_signature(value) != 0) {                   // line 1543
png_icc_tag_name(message+pos, ...);                   // line 1546
png_chunk_benign_error(png_ptr, message);             // line 1566
    \end{lstlisting}

    These lines construct human-readable error messages for invalid profiles or chunk configurations. The execution of these lines indicates that the code was exercised in a stateful, correct parsing context, as triggered by structured inputs. \\ \\

    While \texttt{png.c} showed the most prominent concentration of seed-only coverage, several other source files also exhibited unique coverage only triggered by inputs from the seed corpus:

    \subsubsection*{\texttt{pngerror.c}}
    \begin{lstlisting}
case PNG_NUMBER_FORMAT_x:     // line 178
*--end = digits[number & 0xf]; // line 179
number >>= 4;                  // line 180
    \end{lstlisting}

    These routines are responsible for formatting numerical values in error or debug messages. They happen only when meaningful error conditions occur in the parsing logic, such as during ICC profile evaluation. This type of scenario requires valid PNG inputs and is therefore unlikely to happen without seed corpus.

    \subsubsection*{\texttt{pngmem.c}}
    \begin{lstlisting}
png_error(png_ptr, "Out of memory"); // line 180
    \end{lstlisting}
    This branch occurs in custom memory allocation logic and is reached when parsing  images with large data regions. Structured inputs can induce stressed memory behavior whereas pure mutation struggles to do it.

    \subsubsection*{\texttt{pngrutil.c}}
    \begin{lstlisting}
png_uint_32 profile_length = png_get_uint_32(profile_header); // line 1409
if (png_icc_check_length(...))                                // line 1411
    \end{lstlisting}

    Again, these lines reflect ICC profile length validation logic and chunk-level parsing. The path is inaccessible without an actual iCCP chunk.
    \\ \\
    In all these cases, the lines covered only via corpus-backed fuzzing correspond to:
    \begin{itemize}
        \item \textbf{Structured State Transitions:} Seeded inputs allow the parser to reach and traverse multiple layers of conditionals that random data typically fails to satisfy.
        \item \textbf{Specialized parsing:} Features like ICC chunks parsing are guarded by format flags and values, provided only by a seed corpus.
        \item \textbf{Edge-case error handling:} Some errors occur only with specific inputs that are unlikely to happen under randomness.
        \item \textbf{Memory fallback and numeric formatting paths:} These paths are also only triggerable under valid program states that can be ensured by seeded corpus but are hard to mimic for randomly mutated inputs.
    \end{itemize}


    In conclusion, the comparative evidence confirms that seed corpus inputs do not just enhance coverage breadth, they are critical to reaching code paths that random inputs will have virtually zero chance to explore. Seeded fuzzing allows coverage of realistic application states and error conditions, making it a must for high-quality testing.


    \section{Analyzing Existing Fuzzing Harnesses}

    \texttt{pngread.c} is almost not covered at all, with a line coverage of  3.50\% (26/742). It is responsible The following image metadata can be treated by pngget.c. Those are : \texttt{pHYs, oFFs, READ, bKGD, cHRM, gAMA, sRGB, iCCP, sPLT, cLLI,  mDCV, eXIf, hIST, pCAL, sCAL, sBIT, TEXT, tIME, tRNS}.
    This metadata is critical for many applications. For example, using pHYs metadata, it is possible to specify a conversion between pixels and real world mesurements. It is used by applications requiring accurate physical dimensions (e.g., medical imaging, printing). Uncovered code here implies untested handling of resolution metadata. This is probably an oversight as most PNG images don't use those functionalities, but proper testing is important as those file format are used in professional applications.

    This is probably because the seed has not all the available images metadata, and formats that are handled in the \texttt{pngread.c}. A more comprehensive corpus should fix the issue.

    \section{Improving the Existing Fuzzers}

    \subsection{Improving png\_get}

    To improve the coverage of png\_get, we can add an extra function to the current fuzzing harness that will only    test the png\_get function. This funcion, called \texttt{check\_metadata}, will be called in the fuzzing harness.
    It will take the PNG hnadler created by the fuzzer as input and run the different possible functions to test the png\_get function.
    However, this results in to high a focus those functions, resulting in a low coverage of the other functions.
    Therefore, we only call this function probabilistically, with a 1 in 7 chance of being called. Check metadata follows
    the following pattern:
    \begin{lstlisting}
        void check_metadata(png_structp png_ptr, png_infop info_ptr) {
            png_uint_32 width = png_get_image_width(png_handler.png_ptr, png_handler.info_ptr);
            // other png_get functions to get image properties
            if(PNG_INFO_<FORMAT>) {
                png_get_<FORMAT>(png_handler.png_ptr, png_handler.info_ptr, &<FORMAT_ARGUMENTS>);
            }
        }
    \end{lstlisting}

    To build the fuzzer, the process is the same as the one used in the first part. The list of commands are given in
    part2/run\_w\_update.sh. The current fuzzer is available on our forked repository. The coverage repository
    is available in part2/report. The coverage of is now
    \begin{center}
        \begin{tabular}{ | l | c | c | }
            \hline
            \textbf{Metric} & \textbf{With Corpus} & \textbf{With improved hook and corpus} \\
            \hline
            Lines & 41.6\%  & 41.55\% (5336/12841)\% (6114/13092)\% \\
            Functions &  50.50\% (202/400) &   63.09\% (253/401)  \\
            Regions &  31.39\% (4785/15245) & 36.39\% (5585/15349)  \\
            \hline
        \end{tabular}
        
    \end{center}

\end{document}
